\chapter{Derivazione dei casi di test}
\label{cap:td-derivazione}

Le tabelle di parametri e i \term{test frame} che compaiono nel capitolo
successivo non sono elenchi arbitrari: sono il prodotto di un metodo. Questo
capitolo lo descrive e ne mostra l'applicazione completa a un caso, così che
le tabelle che seguono siano leggibili come risultato e non come premessa.

\section{Il metodo della category partition}

Il metodo di Ostrand e Balcer~\cite{ostrand1988} deriva i casi di test dalla
specifica di una funzione attraverso quattro passi.

\begin{enumerate}
  \item \textbf{Individuare i parametri e le condizioni dell'ambiente} che
    influenzano il comportamento della funzione. I parametri sono gli
    argomenti; le condizioni dell'ambiente sono lo stato del sistema al
    momento della chiamata.
  \item \textbf{Individuare le categorie}: per ciascun parametro, le
    proprietà che ne determinano il trattamento. Per una stringa, tipicamente
    la lunghezza e il formato; per un riferimento, l'esistenza dell'oggetto
    riferito.
  \item \textbf{Partizionare ogni categoria in scelte}, cioè in classi di
    valori trattati allo stesso modo. Ogni scelta è annotata come valida
    (\id{[valido]}) o non valida (\id{[errore]}).
  \item \textbf{Comporre i test frame} combinando una scelta per categoria,
    applicando i vincoli che escludono le combinazioni impossibili o
    ridondanti.
\end{enumerate}

Senza il quarto passo il numero di combinazioni esplode: con sei categorie da
tre scelte si avrebbero 729 casi. I vincoli lo riducono a un numero
gestibile, e sono di due tipi.

\begin{description}[leftmargin=1.4cm,style=nextline]
  \item[\id{[error]}] Una scelta non valida produce un solo test frame, in cui
    tutti gli altri parametri assumono un valore valido. Non ha senso
    combinare due errori: la funzione si arresta al primo, e il secondo non
    verrebbe mai raggiunto.
  \item[\id{[property]} e \id{[if]}] Una scelta può essere selezionata solo se
    un'altra è stata selezionata. Serve a escludere combinazioni impossibili:
    non si può verificare il formato di un'email se l'email è assente.
\end{description}

\section{Convenzioni di identificazione}

\begin{longtable}{@{}L{3.4cm} L{11.0cm}@{}}
\toprule
\textbf{Elemento} & \textbf{Convenzione} \\
\midrule
\endfirsthead
\toprule
\textbf{Elemento} & \textbf{Convenzione} \\
\midrule
\endhead
\bottomrule
\endlastfoot

Scelta &
Sigla del parametro seguita da un numero progressivo. Esempio: \id{LC1} è la
prima scelta della categoria \enquote{lunghezza del commento}. \\

Test frame &
Insieme delle scelte selezionate. Esempio:
\id{\{LC2, VA2, EF1, TU1\}}. \\

Caso di test &
\id{TC\_<area>\_<numero>}, dove \id{<area>} è la sigla del sottosistema
usata nel RAD. Esempio: \id{TC\_FB\_03}. \\

\end{longtable}

\section{Esempio completo: inserimento di un feedback}
\label{sec:td-esempio}

Si applica il metodo all'operazione
\id{creaFeedback(utenteId, tutorialId, dati)} del sottosistema Gestione
feedback, che realizza \id{RF\_FB\_3}.

\subsection{Passo 1 — Parametri e condizioni dell'ambiente}

\begin{longtable}{@{}L{3.4cm} L{4.0cm} L{7.0cm}@{}}
\toprule
\textbf{Elemento} & \textbf{Tipo} & \textbf{Descrizione} \\
\midrule
\endfirsthead
\toprule
\textbf{Elemento} & \textbf{Tipo} & \textbf{Descrizione} \\
\midrule
\endhead
\bottomrule
\endlastfoot

\id{valutazione} & parametro & Voto numerico assegnato al tutorial. \\
\id{commento} & parametro & Testo del giudizio. \\
\id{tutorialId} & parametro & Tutorial oggetto della valutazione. \\
\id{utenteId} & parametro & Autore, ricavato dalla sessione. \\
Feedback preesistente & ambiente &
Presenza di un feedback della stessa coppia utente-tutorial. \\

\end{longtable}

L'ultima riga è una condizione dell'ambiente e non un parametro: non è
passata alla funzione ma ne determina il comportamento, ed è la ragione per
cui il metodo distingue i due concetti.

\subsection{Passo 2 e 3 — Categorie e scelte}

\begin{longtable}{@{}L{3.0cm} L{3.2cm} L{4.4cm} L{2.6cm}@{}}
\toprule
\textbf{Parametro} & \textbf{Categoria} & \textbf{Scelta} & \textbf{Sigla} \\
\midrule
\endfirsthead
\toprule
\textbf{Parametro} & \textbf{Categoria} & \textbf{Scelta} & \textbf{Sigla} \\
\midrule
\endhead
\bottomrule
\endlastfoot

\multirow{4}{*}{\id{valutazione}} & \multirow{2}{*}{Intervallo} &
  $< 1$ \hfill \id{[errore]} & \id{VA1} \\
 & & $1 \le v \le 5$ \hfill \id{[valido]} & \id{VA2} \\
 & & $> 5$ \hfill \id{[errore]} & \id{VA3} \\
 & Tipo & non intero \hfill \id{[errore]} & \id{VA4} \\
\midrule

\multirow{3}{*}{\id{commento}} & \multirow{3}{*}{Lunghezza} &
  $< 2$ caratteri \hfill \id{[errore]} & \id{LC1} \\
 & & da 2 a 500 caratteri \hfill \id{[valido]} & \id{LC2} \\
 & & $> 500$ caratteri \hfill \id{[errore]} & \id{LC3} \\
\midrule

\id{tutorialId} & Esistenza &
  tutorial inesistente \hfill \id{[errore]} & \id{TU1} \\
 & & tutorial esistente \hfill \id{[valido]} & \id{TU2} \\
\midrule

Ambiente & Feedback preesistente &
  presente \hfill \id{[errore]} & \id{FP1} \\
 & & assente \hfill \id{[valido]} & \id{FP2} \\

\end{longtable}

\subsection{Passo 4 — Vincoli e test frame}

I vincoli applicati sono due, entrambi del tipo \id{[error]}:

\begin{itemize}
  \item una scelta non valida di \id{valutazione} o \id{commento} rende
    irrilevanti le scelte successive, perché la validazione dei dati precede
    ogni accesso al database;
  \item la verifica dell'esistenza del tutorial precede quella del feedback
    preesistente, quindi \id{TU1} rende irrilevante la scelta su \id{FP}.
\end{itemize}

Le combinazioni passano così da $4 \times 3 \times 2 \times 2 = 48$ a otto
test frame.

\begin{longtable}{@{}l L{4.4cm} L{3.4cm} L{4.0cm}@{}}
\toprule
\textbf{Caso} & \textbf{Test frame} & \textbf{Esito atteso} & \textbf{Oracolo} \\
\midrule
\endfirsthead
\toprule
\textbf{Caso} & \textbf{Test frame} & \textbf{Esito atteso} & \textbf{Oracolo} \\
\midrule
\endhead
\bottomrule
\endlastfoot

\id{TC\_FB\_01} & \id{\{VA2, LC2, TU2, FP2\}} & Feedback registrato &
  L'oggetto restituito riporta autore, tutorial, voto e commento; il DAO di
  creazione è stato invocato. \\

\id{TC\_FB\_02} & \id{\{VA1, LC2, TU2, FP2\}} & \id{ValidationError} &
  Nessuna scrittura sul database. \\

\id{TC\_FB\_03} & \id{\{VA3, LC2, TU2, FP2\}} & \id{ValidationError} &
  Come sopra. \\

\id{TC\_FB\_04} & \id{\{VA4, LC2, TU2, FP2\}} & \id{ValidationError} &
  Come sopra: una valutazione non intera è rifiutata anche se compresa
  nell'intervallo. \\

\id{TC\_FB\_05} & \id{\{VA2, LC1, TU2, FP2\}} & \id{ValidationError} &
  Come sopra. \\

\id{TC\_FB\_06} & \id{\{VA2, LC3, TU2, FP2\}} & \id{ValidationError} &
  Come sopra. \\

\id{TC\_FB\_07} & \id{\{VA2, LC2, TU1\}} & \id{NotFoundError} &
  L'esistenza del tutorial è verificata prima dell'unicità del feedback. \\

\id{TC\_FB\_08} & \id{\{VA2, LC2, TU2, FP1\}} & \id{ConflictError} &
  Il feedback preesistente non viene sovrascritto. \\

\end{longtable}

\begin{nota}
\textbf{Un caso che il metodo non produce.} La specifica prevede che l'autore
del feedback sia sempre l'utente indicato dal chiamante, mai un valore
proveniente dalla richiesta. Non è una partizione di un parametro: è una
proprietà della funzione. Le proprietà di questo tipo — che il RAD elenca
come requisiti non funzionali — richiedono casi di test scritti a parte, e
sono raccolte in \cref{sec:td-tc-accessi}.
\end{nota}

\subsection{Dal test frame al codice}

Ogni test frame corrisponde a una asserzione eseguibile. Il caso
\id{TC\_FB\_08} è realizzato così:

\begin{lstlisting}[language=Java,caption={Realizzazione di \id{TC\_FB\_08}.},
                   label={lst:td-tc-fb-08}]
it("impedisce due feedback dello stesso utente sullo stesso tutorial",
   async () => {
     feedbackDao.find.mockResolvedValue(unFeedback());   // FP1

     await expect(
       servizio.creaFeedback(1, 1, datiValidi),          // VA2, LC2, TU2
     ).rejects.toBeInstanceOf(ConflictError);
   });
\end{lstlisting}

Il doppio del DAO realizza la condizione dell'ambiente; gli argomenti
realizzano le scelte sui parametri; l'asserzione è l'oracolo.
